\chapter{Baryogenesis, Cosmic Smoothness, and Hubble Friction}

\lectureoverview{Two small dimensionless numbers organize the evening. The first is the surviving matter--antimatter excess, rounded in the course to one part in \(10^8\). Its dynamical explanation requires baryon-number violation, particle--antiparticle asymmetry, and departure from thermal equilibrium: the three Sakharov conditions. The second is the primordial density contrast, of order \(10^{-5}\) at decoupling. Gravity amplifies rather than erases such ripples, so their smallness motivates inflation. We then prepare the mechanics of inflation by deriving the equation of motion of a homogeneous scalar field in an expanding universe and identifying \(3H\dot\phi\) as Hubble friction.}

\section{The Entropy Question Turned Upside Down}

We shall begin inflation tonight, but first we must finish baryogenesis: the creation of the excess of matter over antimatter. The framework is more secure than any detailed numerical model. That is why the subject can be summarized economically even though its central number has not been derived from first principles. \lecturetimestamp{00:00:08--00:00:37}

The historical route into the problem began in 1980, when Bob Wagoner asked Savas Dimopoulos and Leonard Susskind why the entropy of the universe is so large. Put cosmologically, the question was why there are so many photons for every surviving proton. \lecturetimestamp{00:00:40--00:01:27}

At very high temperature, photons, quarks, antiquarks, electrons, and positrons were all abundant and close to thermal equilibrium. Quarks and antiquarks each had number densities of the same broad order as the photon density. That observation already suggests turning the question around: why did annihilation leave any unmatched matter at all? \lecturetimestamp{00:01:27--00:01:49}

The priority of the structural argument belongs to Andrei Sakharov. Dimopoulos and Susskind independently arrived at closely related conditions in 1980, then found Sakharov's 1967 paper while preparing a historical contribution. The conditions are now justly known by Sakharov's name. \lecturetimestamp{00:01:49--00:03:20}

The connection to entropy can now be made explicit. For blackbody radiation, entropy is proportional to photon number up to a fixed factor of order unity. Indeed,
\begin{equation}
s_\gamma=\frac{4\rho_\gamma}{3T},
\qquad
n_\gamma\propto T^3,
\end{equation}
so a comoving thermal photon gas has \(S_\gamma\propto N_\gamma\). The numerical constants are not the issue; the puzzle is the enormous photon-to-proton ratio. \lecturetimestamp{00:04:49--00:05:55}

\begin{classroomqa}
\audiencequestion{How is the photon abundance related to entropy?}
\lecturerresponse{Blackbody radiation is thermal, and its entropy is proportional to its mean photon number. In this context, asking why the cosmic entropy is large is essentially asking why the cosmic microwave background contains so many photons per proton. The proportionality includes numerical constants, but they are irrelevant to the order-of-magnitude puzzle. \lecturetimestamp{00:04:49--00:05:55}}
\end{classroomqa}

As the universe cooled, particles and antiparticles annihilated, while the radiation remained and eventually decoupled. The surprising fact is therefore not that photons survived; it is that annihilation left a tiny unmatched population of matter. \lecturetimestamp{00:05:56--00:07:18}

The useful dimensionless number is therefore
\begin{equation}
\eta_B
\equiv
\frac{N_B-N_{\bar B}}{N_\gamma}
\sim
\frac{N_p-N_{\bar p}}{N_p+N_{\bar p}}
\sim 10^{-8},
\label{eq:c08-baryon-asymmetry}
\end{equation}
where the last value is the rounded order used in the course, not a precision determination. Once an odd-looking number appears, ``it simply began that way'' ceases to be a satisfying stopping point. A quantitative mechanism should explain why the answer is small and why it has a definite sign. \lecturetimestamp{00:07:18--00:08:29}

Calling the excess \emph{matter} rather than \emph{antimatter} is partly conventional: matter is the member of the pair from which we are made. The physical statement is only that one sign dominates throughout the observable region. \lecturetimestamp{00:08:30--00:08:44}

\section{Baryon Number and the First Condition}

Baryon number is normalized so that a proton or neutron carries \(B=1\). In quark language,
\begin{equation}
B=\frac{N_q-N_{\bar q}}{3}.
\label{eq:c08-baryon-number}
\end{equation}
The factor \(1/3\) is essential because each ordinary baryon contains three valence quarks. Other hadrons may carry baryon number, but the unstable ones eventually decay; the net quark excess is the more general bookkeeping device. \lecturetimestamp{00:08:55--00:10:58}

\begin{classroomqa}
\classquestion{Should the net quark number be multiplied or divided by three?}
\lecturerresponse{It must be divided by three. A proton has three quarks and baryon number one, so \(3B=N_q-N_{\bar q}\). The correction on the board is part of the definition, not a change in the physical argument. \lecturetimestamp{00:09:23--00:10:17}}
\end{classroomqa}

Baryon number resembles electric charge only as a possible conserved quantity. Electric charge also sources a long-range Coulomb field; baryon number does not. A proton creates an electric field because it has electric charge, whereas a neutron has baryon number but no net electric charge. Thus conservation and force generation are logically separate properties. \lecturetimestamp{00:11:03--00:12:15}

\begin{classroomqa}
\audiencequestion{What is the relevant difference between baryon number and electric charge?}
\lecturerresponse{Both can serve as bookkeeping charges, but only electric charge is known to source the long-range electromagnetic field. Baryon number does not produce a corresponding Coulomb force. \lecturetimestamp{00:11:03--00:12:15}}
\end{classroomqa}

If \(B\) were exactly conserved, an initial excess would remain frozen forever. Pair creation and annihilation could change \(N_B\) and \(N_{\bar B}\) separately, but only by equal amounts, leaving their difference unchanged. A dynamical theory that begins from a balanced state therefore cannot work unless baryon number can change. \lecturetimestamp{00:12:16--00:13:06; 00:18:06--00:18:43}

The cleanest schematic example is proton decay. Electric charge and energy allow us to imagine
\begin{equation}
p\longrightarrow e^+ +\gamma,
\qquad
\bar p\longrightarrow e^-+\gamma.
\label{eq:c08-decay-pair}
\end{equation}
These channels are bookkeeping models for baryon violation, not predictions of the dominant experimental channel. They take \(B=+1\) or \(-1\) to zero while preserving electric charge. \lecturetimestamp{00:13:06--00:14:34}

\begin{figure}[htbp]
\centering
\begin{tikzpicture}[x=1cm,y=1cm,>=Latex]
\draw[->,thick] (0,1.7)--(2.0,1.7);
\draw[->,thick] (2.0,1.7)--(4.1,2.55);
\draw[->,thick] (2.0,1.7)--(4.1,0.85);
\node[above] at (0.9,1.7) {\(p,\ B=1\)};
\node[right] at (4.1,2.55) {\(\gamma\)};
\node[right] at (4.1,0.85) {\(e^+\)};
\draw[->,thick] (6.0,1.7)--(8.0,1.7);
\draw[->,thick] (8.0,1.7)--(10.1,2.55);
\draw[->,thick] (8.0,1.7)--(10.1,0.85);
\node[above] at (6.9,1.7) {\(\bar p,\ B=-1\)};
\node[right] at (10.1,2.55) {\(\gamma\)};
\node[right] at (10.1,0.85) {\(e^-\)};
\node at (2.0,0.25) {\(B:1\to0\)};
\node at (8.0,0.25) {\(B:-1\to0\)};
\end{tikzpicture}
\caption{Charge-conjugate schematic channels. Either process violates baryon number, but equal rates would create no preferred sign.}
\label{fig:c08-decay-redraw}
\end{figure}

\begin{classroomqa}
\audiencequestion{Could the proton instead decay into a positron and a neutrino?}
\lecturerresponse{That two-fermion final state does not match the fermion parity of the one-fermion initial state. The schematic \(e^++\gamma\) final state contains one fermion and is adequate for the symmetry argument. A realistic proton-decay theory may contain other particles and channels, but the present purpose is only to exhibit \(\Delta B\ne0\). \lecturetimestamp{00:14:39--00:15:07}}
\end{classroomqa}

\begin{classroomqa}
\audiencequestion{Would exact baryon conservation correspond to a symmetry?}
\lecturerresponse{Yes. A conserved quantity may be associated with a continuous phase symmetry, although merely naming that symmetry does not independently prove that nature respects it. The operational test remains whether baryon-number-changing processes occur, so the argument becomes circular if the symmetry is inferred only from the conservation law. \lecturetimestamp{00:15:08--00:15:46}}
\end{classroomqa}

\begin{classroomqa}
\audiencequestion{Can proton--antiproton annihilation itself be regarded as proton decay?}
\lecturerresponse{The terminology is secondary; the conservation law is not. In \(p+\bar p\to\text{photons}\), the initial baryon number is \(1-1=0\), and the final baryon number is also zero. Annihilation removes two particles without changing net \(B\), whereas the one-particle channel in equation~\eqref{eq:c08-decay-pair} changes \(B\). \lecturetimestamp{00:15:48--00:16:47}}
\end{classroomqa}

Present protons are extraordinarily stable: the lecture quotes bounds of order \(10^{32}\)--\(10^{34}\) years, vastly longer than the age of the universe. That observation does not establish exact conservation; it tells us only that baryon-number-changing processes are exceedingly slow at present energies. \lecturetimestamp{00:16:48--00:18:43}

\section{Discrete Symmetries and the Need for a Bias}

The next question is whether the laws distinguish matter from antimatter. Three discrete transformations organize the answer:
\begin{align}
C&:\ \text{particle}\longleftrightarrow\text{antiparticle},\\
P&:\ \mathbf x\longmapsto-\mathbf x,\\
T&:\ t\longmapsto-t.
\end{align}
Here \(P\) reverses spatial handedness, and \(T\) reverses a microscopic movie. The macroscopic second law need not look time-reversal invariant, but the question concerns the underlying microscopic dynamics. \lecturetimestamp{00:18:45--00:22:40}

\begin{figure}[htbp]
\centering
\includegraphics[width=0.82\textwidth]{lecture_08_figure_03.png}
\caption{The board carries the corrected definition \(B=(N_q-N_{\bar q})/3\) into the introduction of \(C\), charge conjugation. The unfinished writing after the word charge is not used as independent evidence.}
\label{fig:c08-charge-board}
\lectureframe{00:20:23.280}
\end{figure}

Relativistic local quantum field theory does not require \(C\), \(P\), or \(T\) separately to be symmetries. Under its standard assumptions, it does require the combined transformation \(CPT\). If a process is possible, replacing every particle by its antiparticle, reflecting the spatial configuration, and reversing the motion gives another allowed process. \lecturetimestamp{00:22:41--00:23:55}

Time reversal in quantum mechanics is subtler than merely changing the sign of \(t\). For a spinless wavefunction in a time-reversal-invariant potential,
\begin{equation}
i\hbar\frac{\partial\psi}{\partial t}=H\psi,
\qquad
T:\ \psi(\mathbf x,t)\longmapsto\psi^*(\mathbf x,-t).
\end{equation}
Changing \(t\to-t\) reverses the time derivative; complex conjugation also sends \(i\to-i\), restoring the equation. For particles with spin, the antiunitary time-reversal operator carries additional structure, but the complex conjugation is already visible in this elementary case. \lecturetimestamp{00:23:58--00:25:04}

\begin{classroomqa}
\audiencequestion{For a wavefunction, does time reversal require complex conjugation as well as \(t\to-t\)?}
\lecturerresponse{Yes. The Schrödinger equation shows why: reversing the time derivative alone changes one side inconsistently, while complex conjugation reverses the sign of \(i\) as well. \lecturetimestamp{00:23:58--00:25:04}}
\end{classroomqa}

We can now state the first Sakharov condition. If the universe began with equal populations of particles and antiparticles, pair creation and annihilation would preserve zero net baryon number unless some process could change \(B\):
\begin{equation}
\boxed{\text{baryon-number violation}.}
\end{equation}
Without it, a symmetric initial state can never acquire a nonzero baryon excess. \lecturetimestamp{00:25:12--00:28:58}

The present stability of the proton is compatible with this condition. In grand-unified reasoning, baryon violation can be mediated by a very heavy particle of mass \(M_X\). At energies far below \(M_X\), the effective amplitude is suppressed schematically by inverse powers,
\begin{equation}
\mathcal A_{\Delta B\ne0}\propto \frac{E^2}{M_X^2},
\qquad
\Gamma_p\sim \frac{m_p^5}{M_X^4}
\quad\text{(dimension-six scaling)},
\label{eq:c08-heavy-suppression}
\end{equation}
up to couplings and channel-dependent factors. A very large \(M_X\) can make a cold proton effectively permanent while allowing baryon-changing reactions in a sufficiently energetic early plasma. \lecturetimestamp{00:29:03--00:33:49}

\begin{classroomqa}
\audiencequestion{Is the extra collision energy already included in the heavy mass \(M_X\)?}
\lecturerresponse{No. \(M_X\) is the mediator's rest-mass scale; \(E\) is the energy available in the collision. The low-energy expansion in \(E/M_X\) is strongly suppressed when \(E\ll M_X\), while the suppression weakens as the thermal collision scale approaches the heavy threshold. The literal denominator written heuristically at the board should be understood in this effective-field-theory sense. \lecturetimestamp{00:32:32--00:34:02}}
\end{classroomqa}

Suppose the initial state contained equal populations of particles and antiparticles. Baryon violation alone would still be unbiased. Every decay \(p\to e^++\gamma\) would have a charge-conjugate partner \(\bar p\to e^-+\gamma\), and the hot plasma would also drive the inverse reactions. Equal microscopic probabilities would produce traffic through baryon number but no average drift. \lecturetimestamp{00:34:50--00:37:22}

\begin{classroomqa}
\audiencequestion{What is the baryon-number change in the antiproton diagram?}
\lecturerresponse{The incoming antiproton has \(B=-1\), while the electron and photon have \(B=0\), so the channel changes baryon number from \(-1\) to zero. Its conjugate changes \(+1\) to zero. Equal occurrence of the two channels gives no net preference. \lecturetimestamp{00:37:24--00:37:39}}
\end{classroomqa}

Random fluctuations are far too small to account for the observed excess. If \(N\) independent events have equal probabilities for two outcomes, then
\begin{equation}
\Delta N\sim\sqrt N,
\qquad
\frac{\Delta N}{N}\sim N^{-1/2}.
\end{equation}
Using the lecture's rough count \(N\sim10^{80}\) for baryons in the observable universe gives
\begin{equation}
\frac{\Delta N}{N}\sim10^{-40},
\end{equation}
thirty-two orders of magnitude below \(10^{-8}\). A coin-toss fluctuation also gives no reason for the same sign to dominate every observed region. \lecturetimestamp{00:37:40--00:40:10}

\begin{classroomqa}
\audiencequestion{Could an initially much larger population, perhaps \(10^{160}\) pairs in the comoving region that became our observable universe, fluctuate by \(10^{80}\)?}
\lecturerresponse{The square-root arithmetic permits that contrived count, but it does not explain the observed spatial coherence. An unrestricted statistical fluctuation would naturally produce matter- and antimatter-dominated patches, not one sign across the visible universe. \lecturetimestamp{00:39:38--00:40:18}}
\end{classroomqa}

The lecture-era observational argument uses two kinds of evidence. Ordinary nuclei, including helium, occur in cosmic rays, whereas antinuclei were not observed in the data being discussed. Antiprotons alone are less decisive because high-energy collisions can create secondary antiprotons. In addition, boundaries between matter and antimatter galaxies would generate conspicuous annihilation radiation. The absence of the required boundary signal strongly disfavors an equal mixture of galaxies and antigalaxies in the observable region. These are observational arguments within the lecture's historical context, not claims about regions beyond our horizon. \lecturetimestamp{00:40:18--00:43:34}

\begin{classroomqa}
\audiencequestion{Are helium nuclei themselves measured in cosmic rays?}
\lecturerresponse{Yes. Charged helium nuclei are measured directly. The detailed cosmic-ray statistics are a separate specialty, but the contrast between common nuclei and the absence of a comparable antinuclear population is the point of the argument. \lecturetimestamp{00:41:53--00:42:22}}
\end{classroomqa}

We therefore need a microscopic bias: violation of \(C\) and, in realistic baryogenesis discussions, \(CP\). Neutral \(B\)-meson systems provide a laboratory example. A neutral meson and its antiparticle can decay to conjugate final states with unequal partial rates,
\begin{equation}
\Gamma(B^0\to f)\ne\Gamma(\bar B^0\to\bar f).
\label{eq:c08-cp-partial}
\end{equation}
The exact channel and numerical asymmetry were not fixed at the board; the robust lesson is the experimentally established existence of particle--antiparticle asymmetry in weak decays. \lecturetimestamp{00:43:38--00:49:13}

The \(CPT\) theorem nevertheless requires equal masses and equal total decay widths for a particle and its antiparticle in vacuum:
\begin{equation}
\Gamma_{\rm tot}(p)=\Gamma_{\rm tot}(\bar p).
\end{equation}
This does not force each branching fraction to agree. Differences in one conjugate pair of channels can be compensated by differences in others, such as schematic electron and muon channels, while the sums remain equal. \lecturetimestamp{00:49:14--00:51:48}

\begin{classroomqa}
\audiencequestion{If proton and antiproton total decay rates are equal, how can a bias arise?}
\lecturerresponse{Equality of total widths does not imply equality channel by channel. Moreover, baryogenesis occurs in a plasma: thermal occupation factors, inverse reactions, and scattering channels enter the kinetic equations. The decisive statement is not that a medium simply evades \(CPT\), but that \(CP\)-asymmetric partial processes can generate a source term once detailed balance is lost. \lecturetimestamp{00:49:39--00:51:48; 00:53:22--00:57:14}}
\end{classroomqa}

This is the second Sakharov condition:
\begin{equation}
\boxed{C\ \text{and}\ CP\ \text{violation}.}
\end{equation}
It supplies directionality to baryon-number-changing reactions. The discovery of \(CP\) violation shortly before Sakharov's paper made this condition physically concrete. \lecturetimestamp{00:51:49--00:53:18}

The minimal Standard Model requires a careful distinction. It does not permit ordinary perturbative proton decay, but its electroweak anomaly allows nonperturbative violation of \(B+L\); at high temperature, sphaleron transitions can make that violation efficient. This is the corrected content of the board's later self-revision. Whether Standard Model sources alone quantitatively reproduce the cosmic asymmetry is a different and much harder question. \lecturetimestamp{00:16:48--00:17:11; 00:52:37--00:53:18}

\section{Equilibrium and the Third Condition}

One obstruction remains. In complete thermal equilibrium, every reaction is balanced by its inverse. If baryon number can change, its equilibrium chemical potential cannot preserve an arbitrary initial \(B\), and detailed balance erases any generated excess. The lecture expresses the same point through \(CPT\): a static equilibrium bath has no macroscopic arrow of time, so the \(T\) factor cannot provide the missing distinction between a process and its reverse. \lecturetimestamp{00:57:14--00:59:42}

The actual universe is expanding and cooling. The macroscopic state therefore distinguishes forward from backward time. This matters only when the expansion competes successfully with microscopic relaxation:
\begin{equation}
\Gamma_{\rm reaction}\lesssim H.
\label{eq:c08-freezeout-condition}
\end{equation}
If \(\Gamma_{\rm reaction}\gg H\), collisions rapidly restore equilibrium despite the expansion. If the reaction rate falls below the Hubble rate, the plasma cannot readjust and the abundance freezes out. \lecturetimestamp{00:59:42--01:02:26}

\begin{classroomqa}
\audiencequestion{Why is expansion alone not automatically enough to take the system out of equilibrium?}
\lecturerresponse{A very slowly expanding bath remains arbitrarily close to equilibrium because microscopic reactions continually readjust it. Departure becomes effective only when the expansion or cooling time \(H^{-1}\) is short compared with the relevant reaction time \(\Gamma^{-1}\). Then a movie of the macroscopic state reveals which direction is forward. \lecturetimestamp{01:01:22--01:02:26}}
\end{classroomqa}

The three Sakharov conditions can now be assembled:
\begin{enumerate}
\item baryon-number violation, so \(B\) is able to change;
\item \(C\) and \(CP\) violation, so the change has a preferred sign;
\item departure from thermal equilibrium, so detailed balance cannot cancel the preference.
\end{enumerate}
They are necessary for the mechanism under discussion and provide a structurally sufficient route to an excess. They do not determine its magnitude without a concrete high-energy model, reaction network, and cosmological history. \lecturetimestamp{01:02:26--01:04:19}

\begin{classroomqa}
\audiencequestion{Does the measured value of the asymmetry constrain high-energy theories?}
\lecturerresponse{In principle, yes. In practice, the result depends on many unknown masses, couplings, phases, reaction rates, and features of the early cosmology. The single measured number is therefore difficult to turn into a narrow allowed region without committing to a detailed model. \lecturetimestamp{01:04:20--01:04:52}}
\end{classroomqa}

\begin{classroomqa}
\audiencequestion{How are the baryon and electron asymmetries connected?}
\lecturerresponse{Electric neutrality links them. In the schematic channel, removing a proton creates a positron; reversing the baryon imbalance shifts the lepton imbalance in the compensating direction. More generally, baryogenesis and leptogenesis are coupled through charge conservation and Standard Model reactions, although the conserved combinations depend on the mechanism. \lecturetimestamp{01:04:54--01:05:59}}
\end{classroomqa}

\begin{classroomqa}
\audiencequestion{Is the observed Standard Model \(CP\) violation too small to explain the baryon excess?}
\lecturerresponse{That is the standard conclusion of detailed electroweak-baryogenesis analyses, but the lecture emphasizes the model dependence of the calculation and the many unknowns at high energy. The defensible statement here is that known low-energy \(CP\) violation exists, while a successful numerical baryogenesis model generally requires additional ingredients. \lecturetimestamp{01:06:01--01:07:50}}
\end{classroomqa}

\begin{classroomqa}
\audiencequestion{Did today's electrons simply come from the destruction of antiprotons?}
\lecturerresponse{The schematic channels make that bookkeeping tempting, but the early plasma contained independent, nearly symmetric lepton populations as well. One may formulate related mechanisms as baryogenesis or leptogenesis, with subsequent reactions transferring asymmetry between sectors. The historical word \emph{baryogenesis} should not be mistaken for a unique microscopic genealogy of every electron. \lecturetimestamp{01:07:51--01:08:44}}
\end{classroomqa}

\begin{classroomqa}
\audiencequestion{Why is the discussion usually framed in terms of baryons rather than all leptons?}
\lecturerresponse{Baryonic matter is inventoried through atomic spectra and other astronomical observations, whereas the cosmic neutrino sector is much harder to count directly. The naming reflects what was observationally accessible as well as the history of the subject. \lecturetimestamp{01:08:45--01:09:12}}
\end{classroomqa}

\section{Annihilation Leaves the Excess Behind}

At high temperature, pair creation and annihilation run in both directions,
\begin{equation}
\gamma+\gamma\longleftrightarrow p+\bar p,
\qquad
E_{\gamma\gamma}\gtrsim2m_pc^2\approx2\,\mathrm{GeV}.
\end{equation}
As the bath cools, photons energetic enough to create a baryon pair become exponentially rare. Annihilation remains energetically allowed, but inverse production can no longer maintain equilibrium. Almost every proton finds an antiproton, and the small excess in equation~\eqref{eq:c08-baryon-asymmetry} is what remains. The energetic annihilation photons quickly share their energy with the thermal bath. \lecturetimestamp{01:09:13--01:11:44}

\begin{classroomqa}
\audiencequestion{Is the ratio \(10^{-8}\) a present number, and when did the large annihilation occur?}
\lecturerresponse{It is the course's present baryon-to-photon order of magnitude. Earlier, both baryons and antibaryons were numerous. Once inverse pair creation became ineffective, their symmetric component annihilated and left the net excess. \lecturetimestamp{01:09:13--01:11:44}}
\end{classroomqa}

\begin{classroomqa}
\audiencequestion{When did proton--antiproton annihilation stop?}
\lecturerresponse{It stopped when antiprotons had been depleted to a negligible abundance. The lecture locates the event only broadly: at GeV-scale temperatures, well before nucleosynthesis and certainly long before recombination. No precise time is inferred from the discussion. \lecturetimestamp{01:11:45--01:12:13}}
\end{classroomqa}

The annihilation photons have not disappeared. They belong to the thermal radiation history that ultimately feeds the microwave background. Radiation nevertheless contributes little to the energy budget today because
\begin{equation}
\rho_{\rm rad}\propto a^{-4},
\qquad
\rho_m\propto a^{-3}.
\end{equation}
Photon number is large, but each photon's energy redshifts; matter retains its rest-mass energy. \lecturetimestamp{01:12:14--01:13:18}

\begin{classroomqa}
\audiencequestion{If annihilation made so many photons, why is radiation dynamically negligible today?}
\lecturerresponse{The photons are still present, chiefly as the microwave background, but their energy density loses an extra power of the scale factor. Expansion both dilutes their number and redshifts each photon. \lecturetimestamp{01:12:14--01:13:18}}
\end{classroomqa}

\begin{classroomqa}
\audiencequestion{Does a very heavy mediator mean the Standard Model is incomplete?}
\lecturerresponse{A grand-unified proton-decay mediator would indeed be new physics beyond the minimal Standard Model. The later correction about anomalous Standard Model baryon violation does not supply such an ordinary proton-decay particle, nor by itself a complete quantitative baryogenesis model. \lecturetimestamp{01:13:20--01:13:31}}
\end{classroomqa}

\begin{classroomqa}
\audiencequestion{Could a large inequality between the two schematic decay channels convert an initially pure \(p\)-\(\bar p\) population into the observed equal electron and proton counts?}
\lecturerresponse{One can invent such bookkeeping, but a hot equilibrium plasma already contains large populations of all accessible particles and antiparticles. Moreover, observed \(CP\) asymmetries are not a license to make one channel arbitrarily dominant. A consistent model must evolve the full coupled reaction network rather than only the two arrows on the board. \lecturetimestamp{01:13:33--01:15:10}}
\end{classroomqa}

This is where the baryogenesis discussion honestly ends. The qualitative ingredients are known; the observed number is not derived. A more complete answer awaits both a more complete particle theory and a controlled history of the very early universe. \lecturetimestamp{01:15:13--01:15:28}

\section{The Second Number: Primordial Lumpiness}

Inflation addresses another question that became quantitative through the cosmic microwave background: why is the universe so homogeneous and isotropic on large scales? The CMB has an extraordinarily precise blackbody spectrum with temperature near \(2.7\,\mathrm K\), and that temperature is nearly the same in every direction. What had once been a broad cosmological impression became a high-precision fact. \lecturetimestamp{01:15:28--01:17:27}

Gravity makes the smoothness more surprising. Diffusion spreads a drop of ink until its concentration becomes uniform. Attractive gravity does the opposite: an overdense region pulls in surrounding matter, increasing its density while emptying neighboring regions. Small perturbations therefore seed galaxies, clusters, and voids by growing with time. \lecturetimestamp{01:17:28--01:22:12}

\begin{figure}[htbp]
\centering
\begin{tikzpicture}[x=1cm,y=1cm,>=Latex]
\draw[thick] (-4.2,0) ellipse (1.25 and 0.72);
\draw[very thick] (0,0) ellipse (1.55 and 0.92);
\draw[thick] (4.2,0) ellipse (1.25 and 0.72);
\node at (-4.2,0) {underdense};
\node at (0,0) {overdense};
\node at (4.2,0) {underdense};
\draw[->,thick] (-2.65,0)--(-1.65,0);
\draw[->,thick] (2.65,0)--(1.65,0);
\draw[->,thick] (-2.85,0.75)--(-1.45,0.4);
\draw[->,thick] (2.85,0.75)--(1.45,0.4);
\draw[->,thick] (-2.85,-0.75)--(-1.45,-0.4);
\draw[->,thick] (2.85,-0.75)--(1.45,-0.4);
\end{tikzpicture}
\caption{Gravitational instability transfers matter from underdense surroundings into an overdensity, amplifying an initial contrast rather than diffusing it away.}
\label{fig:c08-gravity-growth}
\end{figure}

Running this growth backward implies a much smoother early universe. Exact homogeneity would also be unacceptable: exact symmetry would persist and no galaxies could form. The useful measure is the density contrast
\begin{equation}
\delta\equiv\frac{\delta\rho}{\rho},
\end{equation}
where \(\delta\rho\) is the local departure from the background density. Structure-growth estimates associated with Jim Peebles and other early cosmologists required a primordial amplitude around
\begin{equation}
\boxed{\frac{\delta\rho}{\rho}\sim10^{-5}}
\end{equation}
at decoupling, with \(10^{-4}\)--\(10^{-5}\) quoted as the historical range. Once again, the number sharpens the question without yet explaining its value. \lecturetimestamp{01:21:20--01:24:17}

The broad inflationary proposal, introduced by Alan Guth in 1980, is a period of enormous, approximately exponential expansion. Stretching a small smooth region by many orders of magnitude dilutes pre-existing curvature and pushes irregularities to scales larger than the observable patch. The crinkled-balloon analogy captures the smoothing intuition, though it should not be taken as a literal embedding of space. \lecturetimestamp{01:24:18--01:25:35}

\begin{classroomqa}
\audiencequestion{How can the universe be isotropic when nearby matter is visibly lumpy in different directions?}
\lecturerresponse{Isotropy is a large-scale statistical statement, not a claim that this room, the Galaxy, or every line of sight looks identical. Earth is mountainous locally yet extremely smooth compared with its radius. Likewise, galaxies and clusters are local structure superposed on a very smooth large-scale background. \lecturetimestamp{01:25:37--01:26:22}}
\end{classroomqa}

\begin{classroomqa}
\audiencequestion{Wouldn't a random process occasionally create enormously overdense directions?}
\lecturerresponse{The perturbations are not unrestricted independent fluctuations. They form a random field with a particular, nearly homogeneous and isotropic statistical spectrum. Their amplitudes are small, of order \(10^{-5}\), even when their wavelengths are large. The spectrum and its origin are the subjects to be developed, not assumptions of arbitrarily large coin-toss excursions. \lecturetimestamp{01:26:26--01:28:23}}
\end{classroomqa}

\begin{classroomqa}
\audiencequestion{Was inflation caused by the high temperature associated with early times?}
\lecturerresponse{Temperature is not the main variable during inflation. The central phenomenon is the rapid exponential expansion driven by the stress-energy of a field. A thermal bath becomes important again when inflation ends and its energy is converted into particles. \lecturetimestamp{01:28:25--01:29:04}}
\end{classroomqa}

\section{A Mechanical Prelude: Motion with Drag}

To prepare the inflation equations, consider a unit-mass stone falling through a viscous fluid. Call its vertical coordinate \(\phi\), anticipating the scalar field that will soon replace the stone. If its potential energy is \(V(\phi)\), the conservative force is
\begin{equation}
F(\phi)=-\frac{dV}{d\phi},
\label{eq:c08-force-potential}
\end{equation}
and without drag Newton's equation is \(\ddot\phi=F\). \lecturetimestamp{01:29:39--01:32:58}

\begin{figure}[htbp]
\centering
\includegraphics[width=0.82\textwidth]{lecture_08_figure_06.png}
\caption{The readable board relation \(F=-dV/d\phi\) for the falling-stone analogy. The downward arrow fixes the convention that increasing \(\phi\) points upward.}
\label{fig:c08-force-board}
\lectureframe{01:35:09.513}
\end{figure}

For linear viscous drag, the additional force is \(-\gamma\dot\phi\), with \(\gamma>0\). The equation becomes
\begin{equation}
\ddot\phi+\gamma\dot\phi=F(\phi)
=-V'(\phi).
\label{eq:c08-viscous-motion}
\end{equation}
Drag vanishes at rest and always opposes the velocity. If \(\phi\) increases upward, a falling stone has \(F<0\) and \(\dot\phi<0\); then \(-\gamma\dot\phi>0\), so drag points upward as required. \lecturetimestamp{01:32:58--01:35:27}

\begin{classroomqa}
\classquestion{Why do the potential force and the drag force appear with the same minus sign? Should they not point in opposite directions?}
\lecturerresponse{Their directions are determined by different quantities. The potential force is \(-V'(\phi)\). Drag is \(-\gamma\dot\phi\). During downward motion, both \(F\) and \(\dot\phi\) are negative, so \(F\) points down while \(-\gamma\dot\phi\) points up. The algebra therefore produces opposite physical directions without changing either formula. \lecturetimestamp{01:34:16--01:35:27}}
\end{classroomqa}

As the stone accelerates, the drag grows until the forces balance. For approximately constant \(F\), terminal motion satisfies
\begin{equation}
0=F-\gamma\dot\phi_{\rm term},
\qquad
\dot\phi_{\rm term}=\frac{F}{\gamma}.
\end{equation}
The sign is correct: a downward \(F\) gives a negative terminal velocity. A shallow potential and large \(\gamma\) therefore produce slow motion down the hill. \lecturetimestamp{01:35:27--01:38:18}

\begin{figure}[htbp]
\centering
\begin{tikzpicture}[x=1cm,y=1cm,>=Latex]
\draw[->] (0,0)--(7.0,0) node[right] {\(\phi\)};
\draw[->] (0,0)--(0,4.4) node[above] {\(V(\phi)\)};
\draw[very thick] plot[smooth] coordinates {(0.5,3.7) (1.6,3.55) (2.8,3.2) (4.0,2.55) (5.1,1.65) (6.3,0.6)};
\filldraw (1.45,3.58) circle (0.09);
\draw[->,thick] (1.55,3.48)--(2.25,3.23);
\node[align=left,anchor=west] at (2.45,3.75) {large drag,\\slow downhill motion};
\end{tikzpicture}
\caption{A shallow potential with strong drag. The field analogue will replace the stone's coordinate by a spatially homogeneous scalar amplitude.}
\label{fig:c08-shallow-potential}
\end{figure}

\section{The Inflaton and Hubble Friction}

Inflation posits an additional scalar field \(\phi\), conventionally called the inflaton. Historically this was a bold hypothesis introduced to explain cosmic smoothness. We begin with a field approximately uniform in space, \(\phi=\phi(t)\). Expansion stretches spatial variations, while the homogeneous problem exposes the essential dynamics. \lecturetimestamp{01:38:38--01:40:14}

For a canonical scalar field, the full local energy density is
\begin{equation}
\rho_\phi
=
\frac12\dot\phi^{\,2}
+\frac{1}{2a^2}(\boldsymbol\nabla\phi)^2
+V(\phi).
\end{equation}
Under the homogeneous approximation \(\boldsymbol\nabla\phi=0\), this reduces to
\begin{equation}
\rho_\phi=\frac12\dot\phi^{\,2}+V(\phi).
\label{eq:c08-scalar-density}
\end{equation}
The first term is kinetic energy in field configuration space: it measures temporal change, not translational motion through ordinary space. \lecturetimestamp{01:40:18--01:42:20}

\begin{classroomqa}
\audiencequestion{Does the dot on \(\phi\) mean a partial derivative with respect to time?}
\lecturerresponse{Yes. For the homogeneous field, \(\dot\phi=d\phi/dt\). Spatial gradients would contribute a separate positive term, but they vanish under the present assumption. \lecturetimestamp{01:41:08--01:41:34}}
\end{classroomqa}

Follow a unit comoving coordinate box. Its physical volume is \(a^3(t)\), so the scalar energy inside it is
\begin{equation}
E_\phi
=a^3(t)
\left[
\frac12\dot\phi^{\,2}+V(\phi)
\right].
\end{equation}
The explicit factor \(a^3(t)\) is the new ingredient: the physical box expands even though its comoving coordinates are fixed. \lecturetimestamp{01:42:21--01:44:09}

\begin{classroomqa}
\audiencequestion{Does the scale factor \(a\) depend on spatial position?}
\lecturerresponse{Not in the homogeneous FRW model. It is the same throughout space at a given cosmic time but may vary with time. If \(\phi\) depended on position, its spatial-gradient energy would have to be restored. \lecturetimestamp{01:44:12--01:44:49}}
\end{classroomqa}

\begin{classroomqa}
\audiencequestion{Is the field value associated with each point in space even though we now write only one \(\phi(t)\)?}
\lecturerresponse{Yes. A field assigns a value to every spacetime point. Homogeneity means those spatial values agree at each time, reducing the dynamics to one common function \(\phi(t)\). \lecturetimestamp{01:44:50--01:45:06}}
\end{classroomqa}

\begin{classroomqa}
\audiencequestion{What physical motion does the field's kinetic term describe? Is the field becoming stronger or weaker?}
\lecturerresponse{That is determined by the equation of motion, the initial value and velocity, and the sign of \(V'(\phi)\). The analogy is a particle whose direction of motion cannot be known from its energy formula alone. \lecturetimestamp{01:45:07--01:46:08}}
\end{classroomqa}

The homogeneous Lagrangian is kinetic minus potential energy,
\begin{equation}
L_\phi=a^3(t)
\left[
\frac12\dot\phi^{\,2}-V(\phi)
\right].
\label{eq:c08-scalar-lagrangian}
\end{equation}
Applying the Euler--Lagrange equation gives
\begin{align}
\frac{d}{dt}\left(\frac{\partial L_\phi}{\partial\dot\phi}\right)
&=
\frac{\partial L_\phi}{\partial\phi},\\
\frac{d}{dt}\left(a^3\dot\phi\right)
&=-a^3V'(\phi).
\end{align}
The product rule supplies two terms,
\begin{equation}
a^3\ddot\phi+3a^2\dot a\,\dot\phi
=-a^3V'(\phi).
\end{equation}
Dividing by \(a^3\) and defining \(H=\dot a/a\), we obtain
\begin{equation}
\boxed{
\ddot\phi+3H\dot\phi+V'(\phi)=0
}
\label{eq:c08-hubble-friction}
\end{equation}
or, with \(F=-V'(\phi)\),
\begin{equation}
\ddot\phi+3H\dot\phi=F.
\end{equation}
Every line follows from the expanding volume factor in equation~\eqref{eq:c08-scalar-lagrangian}. \lecturetimestamp{01:46:10--01:51:23}

\begin{classroomqa}
\audiencequestion{Is \(H=\dot a/a\) the same everywhere, and is it really a constant?}
\lecturerresponse{In the homogeneous model it is the same at every spatial point on a cosmic-time slice. It need not be constant in time; ``Hubble constant'' often names its present value, while \(H(t)\) is the Hubble rate throughout the evolution. \lecturetimestamp{01:50:32--01:50:57}}
\end{classroomqa}

Equation~\eqref{eq:c08-hubble-friction} is the payoff. It has the same form as viscous motion, with drag coefficient
\begin{equation}
\gamma_{\rm cosmic}=3H.
\end{equation}
Expansion therefore slows the evolution of a homogeneous scalar field. If \(H\) is large and \(V(\phi)\) is sufficiently flat, the inertial term becomes small compared with the drag term and the field approaches the slow-roll relation
\begin{equation}
3H\dot\phi\simeq-V'(\phi).
\end{equation}
Like a ball descending through cold, viscous oil, the field can take an extraordinarily long time to cross a shallow potential even though a force acts on it. \lecturetimestamp{01:51:23--01:53:30}

The next step is to insert the scalar energy density into the Friedmann equation and let the expansion rate respond to the slowly evolving field. That coupled system produces inflation. For now, the essential preparation is complete: the field stores kinetic and potential energy, the expanding volume generates Hubble friction, and a slowly rolling field can keep its potential energy nearly constant long enough to drive rapid expansion. \lecturetimestamp{01:53:30--01:54:11}

\section{What to Carry Forward}

The baryon asymmetry and the primordial density contrast play parallel roles. Neither \(10^{-8}\) nor \(10^{-5}\) is explained merely by naming it, but each converts a vague cosmological impression into a quantitative problem. The first leads to the Sakharov conditions; the second leads toward inflation.

The bridge into inflation is equally precise. A homogeneous scalar field in an FRW universe obeys
\[
\ddot\phi+3H\dot\phi+V'(\phi)=0.
\]
The middle term is not an inserted phenomenological viscosity. It follows from the growth of the physical volume \(a^3(t)\). That single result will organize the next lecture's account of a universe driven by a field rolling slowly down its potential.
